\section{Experiments}

\subsection{Protocol}

We freeze one method configuration across scenes. Mapping images are used for
Gaussian-prior construction and all offline localization-map learning; official
test images are used only for evaluation. The deployment frontend uses native
SuperPoint, full-resolution coordinates with the explicit pixel-center
convention, global cosine top-1 matching, and one PoseLib PnP/RANSAC solve.

We report median, mean, and 90th-percentile translation and rotation error;
recall at $5\,$cm/$5^\circ$; raw and RANSAC-inlier correspondence precision at
2 and 4 pixels; catastrophic failures; primitive, track, and final anchor
counts; map memory; build time; and localization P50/P90 latency.

\subsection{Datasets and Baselines}

The evaluation covers five Cambridge Landmarks scenes, all seven 7Scenes
scenes, and all twelve 12Scenes scenes. The indoor main tables use the same
off-the-shelf Gaussian-prior method without semantic masks. A representative
prior-flexibility subset compares vanilla 3DGS, vanilla 2DGS, an enhanced 2DGS
prior, and a feed-forward AnySplat prior while keeping every localization-map
stage fixed.

Our principal comparison is between the Gaussian-supported KCS/GWFF bootstrap
(A0) and the complete reconstructed map (A1). Track-only and base-only maps are
reported on representative scenes to isolate track identity and Gaussian
coverage.

\begin{table*}[t]
  \centering
  \caption{Frozen-protocol localization results. Values are populated only by
  the registered experiment summaries.}
  \label{tab:main-results}
  \resizebox{\textwidth}{!}{%
  \begin{tabular}{llrrrrrrrr}
    \toprule
    Dataset & Method & Med. TE & Mean TE & P90 TE & Med. AE & 5/5 & Raw P@2 & Anchors & P90 time \\
    \midrule
    Cambridge & A0 & \multicolumn{8}{c}{\TODO{aggregate five-scene results}} \\
    Cambridge & A1 & \multicolumn{8}{c}{\TODO{aggregate five-scene results}} \\
    7Scenes & A0 & \multicolumn{8}{c}{\TODO{run seven official scenes}} \\
    7Scenes & A1 & \multicolumn{8}{c}{\TODO{run seven official scenes}} \\
    12Scenes & A0 & \multicolumn{8}{c}{\TODO{run twelve official scenes}} \\
    12Scenes & A1 & \multicolumn{8}{c}{\TODO{run twelve official scenes}} \\
    \bottomrule
  \end{tabular}}
\end{table*}


\subsection{Analysis Plan}

The evaluation tests four questions: whether A1 consistently improves A0;
whether compact topology preserves both central and tail pose accuracy; whether
the reconstructed map is robust to off-the-shelf prior quality; and whether the
offline cost yields a simpler online pipeline. A scene failure is attributed
only after auditing dataset convention, prior reconstruction, track coverage,
and candidate precision under the frozen protocol.
